\section{数据结构源码视角和容器选择}

数据结构学习要同时看抽象接口和存储模型。vector 的接口像可变长数组，底层却有大小、容量和连续内存三件事；deque 的接口像两端都能高效进出的序列，底层通常是分段连续块；list 的接口保证节点稳定和常数移动，但牺牲随机访问和缓存局部性；unordered\_map 的接口给平均常数查询，底层却受哈希函数、桶数量、负载因子和 rehash 影响。理解这些模型后，刷题时选择容器才不是背名字。

vector 的扩容是最值得做的实验。连续 push\_back 时，size 每次加一，capacity 只在容量不足时跳变，data 地址也可能改变。这个现象解释了为什么保存 vector 元素指针或迭代器后继续插入很危险，也解释了 reserve 和 resize 的区别。reserve 改容量，不构造新元素；resize 改大小，会构造或销毁元素。很多性能问题不是算法复杂度错，而是无意中触发了大量重新分配和拷贝。

string 的成本不能只看语法。substr 往往会构造新字符串，动态规划中两层循环里反复 substr 可能把复杂度悄悄提高一阶。string\_view 能表达非拥有视图，但它要求底层字符串生命周期足够长。刷题中为了清晰可以先用 substr，复盘时必须知道它的复制成本；工程代码中，如果热路径上大量切片，应该考虑下标范围、视图或 Trie。

unordered\_map 的平均常数时间不是无条件承诺。哈希冲突、负载因子过高、频繁 rehash 都会让性能波动。刷题时调用 reserve 可以减少 rehash，使用 find 可以避免 operator[] 的默认插入副作用。当前缀和计数题里，哈希表语义是历史前缀频次，如果用 operator[] 随意读取不存在键，就可能把调试输出和真实状态混在一起。

map 和 set 基于有序结构，常用于前驱、后继、区间查询和动态有序集合。它们的查询是对数复杂度，不如哈希表平均常数快，但能回答哈希表回答不了的顺序问题。区间覆盖、日程安排、滑动窗口有序统计、离线扫描线，都可能需要有序容器。选择 map 不是因为它高级，而是因为题目需要顺序语义。

priority\_queue 是容器适配器，不是完整排序容器。它只承诺堆顶极值，不承诺遍历顺序。需要频繁删除任意元素时，普通 priority\_queue 不支持直接删除，常见办法是懒删除：另一个哈希表记录应该删除的元素，真正取堆顶前不断弹出过期项。理解这个限制，才能写好滑动窗口中位数、任务调度和带过期时间的事件模拟。

Linux 内核的数据结构会在独立专题中系统展开。这里先只保留一个接口判断：STL 容器强调泛型接口、异常安全和所有权封装；内核结构强调嵌入式节点、明确生命周期、分配控制和并发语义。两者都不是模板清单，核心仍然是访问模式。

读源码时不要从模板细节硬啃。先读接口语义和复杂度承诺，再找核心成员，最后追一条热路径。vector 看 push\_back 和扩容，deque 看段表如何定位下标，unordered\_map 看查找和 rehash，map 看 lower\_bound 如何沿树下降，priority\_queue 看 push\_heap 和 pop\_heap。每次阅读只回答一个问题：这个 API 为什么是这个复杂度，什么时候会让迭代器失效，题目里用它是否合适。

\begin{principlebox}{专题复盘要求}
读完本节后，不要只记结论。请把每个结论还原成一个可验证的问题：它解决了哪一步重复工作，依赖哪个题目条件，使用哪个容器或状态，边界条件如何破坏错误写法。
\end{principlebox}
